\documentclass[11pt]{article}

\usepackage[margin=1in]{geometry}
\usepackage{amsmath,amssymb,amsthm}
\usepackage{enumitem}
\usepackage{microtype}
\usepackage[hidelinks]{hyperref}
\usepackage[nameinlink,noabbrev]{cleveref}
\setlength{\emergencystretch}{4em}

\newcommand{\PP}{\mathbb{P}}
\newcommand{\EE}{\mathbb{E}}
\newcommand{\one}{\mathbf{1}}

\newtheorem{definition}{Definition}
\newtheorem{remark}{Remark}

\title{Threat Models, Observation Channels, and Repetition Windows for Anonymous DHTs:\ A Short Declaration Checklist for Referee-Grade Budget Claims}
\author{}
\date{June 18, 2026 (archive rev0900)}

\begin{document}
\maketitle

\begin{abstract}
Anonymous DHT papers are easy to read past each other because ``the threat model'' hides several independent choices:
what the secret is (initiator vs target), what transcript projection is visible (tier), what the repetition unit/window is, and how much observation/manipulation power the adversary has.
This short note standardizes a small declaration checklist meant to be cited (not re-derived) by the Boss-Fight dial sheet, certified-menu selection, and evaluation/notarization papers.
The central correction in rev0900 is that an observation rate, compromise fraction, or committee-hit probability is telemetry, not an observation channel. A budget-bearing declaration must identify the complete conditional transcript law or a source-bound theorem/witness that dominates it, including adaptive history, retries, stopping, aborts, and visible length.
It is complementary to the series trace-interface note, which owns the transcript-projection vocabulary.
\end{abstract}


\paragraph{A short ``TW declaration'' template.}
At minimum, a referee should be able to locate these fields: secret (\emph{initiator} vs \emph{destination key}), observation tier (by \texttt{exposure\_nf\_id}), repetition unit, window, adversary control (passive vs selective-failure vs full Byzantine), and background conditioning (public calendars, committee partitions).
This note treats that tuple as the named object \(\mathsf{TW}\) whose digest is \texttt{tw\_id}.


\begin{remark}[A concrete template instance]
Synthesis~17's worked example instantiates this template with the label
\texttt{tw.reader\_24h.q500.v1} and a digest-bound \texttt{tw\_id} that is reused across multiple receipt line items.
The point is mechanical: a referee should be able to diff the window declaration (and its conditioning universe) without reading prose.\cite{series:workedexample,series:receiptitems}
\end{remark}


\paragraph{Scope.}
Standardize threat-model and repetition-window declarations for Anonymous-DHT anonymity budgets (so papers can stay non-redundant).

\paragraph{Imports.}
We import transcript-projection vocabulary (observation tiers / exposure normal forms) from the trace-interface synthesis note.\cite{series:traceifaces}
Receipt line items bind these declarations via digest IDs (\texttt{exposure\_nf\_id}, \texttt{tw\_id}); see Synthesis~12.\cite{series:receiptitems}
Primary-path role-separation assumptions (and their optional probabilistic adjuncts) are imported from the separation-manifest notes.\cite{series:primarysepmanifest,series:probsep}

\paragraph{Exports.}
A named tuple for secret scope, visible tier, repetition unit/window, conditioning universe, observation channel, and adversary powers.

\paragraph{Artifact policy.}
Source-only. A channel-dependent claim is publication-ineligible unless its witness and checker are present in the bound release packet.

\section{What this note owns (non-redundancy)}
The series trace-interface note owns \emph{what is visible} (observation tiers, exposure normal forms, and the CCT schema)~\cite{series:traceifaces}.
This note owns \emph{how long} and \emph{under what repetition}:
(i) the repetition unit/window declaration; (ii) the candidate universe and conditioning rules that define the secret scope; and (iii) a minimal adversary-power menu (passive, active/selective-failure, background-manipulation).
Other papers should cite this note instead of repeating generic ``Boss Fight vs passive'' paragraphs.

\section{A minimal declaration checklist}\label{sec:checklist}
Any referee-facing anonymity claim in this archive should explicitly declare the following fields.
(These fields are also natural inputs to ABOM/OINL certificates.)

\begin{definition}[Threat/window declaration]
A threat/window declaration is a tuple
\[
\mathsf{TW}=(\mathsf{secret},\ \mathsf{tier},\ \mathsf{unit},\ \mathsf{window},\ \mathsf{universe},\ \mathsf{obs},\ \mathsf{active},\ \mathsf{bg}).
\]
The fields are:
\begin{enumerate}[leftmargin=*]
\item \textbf{Secret.} What is being protected: initiator identity $I$, lookup target/key $K$, destination committee label $A\in\mathcal A$, or a profile/class $C$.
\item \textbf{Tier.} Which transcript projection is being bounded (Tier~0--3) and its exposure normal form~\cite{series:traceifaces}.
\item \textbf{Unit.} The repetition unit for composition (per lookup, per epoch, per round-trip, per retry attempt).
\item \textbf{Window.} The window length/horizon for the claimed budget (e.g., $T$ epochs, or ``per 24h sliding window'') and the reset/checkpoint rule.
\item \textbf{Universe.} The candidate universe size and conditioning: e.g., $|\mathcal A|$ committees at Tier~0, or an effective target set size $K$ after public side information, caching, and eligibility disclosures.
\item \textbf{Observation channel.} The complete conditional law of the declared transcript given the secret and reachable pre-release history, or a named channel family with a source-bound witness and checker proving every precondition. Scalar observation or compromise rates are telemetry, not substitutes for the channel. Include retries, stopping/abort decisions, visible length, path selection, and public side information used in conditioning.
\item \textbf{Active power.} What deviations are in-scope: selective failures, induced retries/delays, Sybil/route capture, equivocation in logged receipts.
\item \textbf{Background assumptions.} If the claim references a public background $Q$ (equalization), declare how $Q$ is produced (specified vs measured), how drift is detected, and what manipulation is in-scope.
\end{enumerate}
\end{definition}

\paragraph{Deployed control-plane.}
If $\mathsf{obs}$ or $\mathsf{bg}$ are defined in terms of delegated-routing endpoints, AutoConf, gateway defaults, or other IPFS/libp2p control-plane knobs,
treat those knobs as part of the declared state and see Synthesis~30 for a compact ``what-to-pin'' checklist.\cite{series:deployedsurfaces}

\begin{definition}[Threat-window id (TW-ID)]
Like interfaces, threat/window declarations should be treated as claim fields.
Define a stable identifier
\[
\mathsf{tw\_id}=\mathrm{SHA256}(\mathrm{canon}(\mathsf{TW}))\in\{0,1\}^{256},
\]
where $\mathrm{canon}(\cdot)$ is a canonical serialization of the eight named fields in $\mathsf{TW}$, including the exposure-normal-form id (ENF-ID), channel-witness type and digest, adaptive-history convention, and visible stop/abort/length behavior when applicable.
Publish $\mathsf{tw\_id}$ as \texttt{sha256:<lowercase-hex>} (explicit algorithm tag). A short prefix may be used in text provided the full digest is in the ABOM.
\end{definition}

\begin{remark}[Field naming (avoid case drift)]
In this archive, the claim field is written as \texttt{tw\_id} (lowercase) in receipts, ABOMs, and OINLs.
Mixing \texttt{TW\_id} and \texttt{tw\_id} as distinct keys is an avoidable source of audit drift; treat the id as a single field.
\end{remark}

\begin{remark}[Why TW-ID matters]
Without TW-ID, a certificate can silently change its composition horizon (``per epoch'' vs ``per day''), its candidate universe, or its permitted active deviations.
TW-ID makes these choices explicit and supports comparison across releases and across papers.
\end{remark}

\begin{remark}[Why these fields are separable]
A Tier~0 ``single-label'' guarantee can be meaningful even if Tier~2 is hopeless, but only if the deployment enforces interface discipline (logs/certificates do not accidentally publish higher-tier traces).
Similarly, a strong per-lookup bound is meaningless if the deployment allows adversarially induced retries that expand the repetition window.
Splitting the declaration into \emph{tier} and \emph{window} avoids silent mismatches.
\end{remark}

\begin{remark}[Separation assumptions belong in the declaration]
If a claim relies on a ``primary path'' whose security is conditional on non-collusion (e.g., client\,$\to$\,relay\,$\to$\,gateway), the separation assumption is part of the threat model, not an implementation footnote.
In this series, treat it as a digest-bound guard object (a separation-manifest id) and, when publishing a quantitative failure bound, as a guarded checkpoint via \texttt{rho\_upper} (Synthesis~23 gives the receipt-side $(0,\rho)$ interpretation). \cite{series:primarysepmanifest,series:probsep}
\end{remark}

\section{Two threat-model templates (used across the archive)}\label{sec:templates}
We intentionally keep templates few.
Papers should specify deviations from these templates rather than introducing new ones casually.

\paragraph{Template P (passive profiling window).}
Secret is a class/identity variable $C$.
Tier is 0 or 1.
Unit is ``epoch'' (or lookup).
Window is long ($T\gg 1$).
The adversary is passive and observes the declared tier projection each unit.
Background $Q$ may drift and must be logged/budgeted if it is used in a contract.
This template is the default for the profiling/equalization paper (Anonymity~B).

\paragraph{Template B (Boss Fight window).}
Secret is a target/destination $K$ (or committee label) and the adversary can actively shape the transcript: selective failures, induced retries, and timing perturbations.
Tier is typically 2 (ordered failures) and the repetition unit is ``attempt'' or ``lookup''.
The window is defined by an operational query budget $Q$ (lookups per time window) together with a reset rule.
This template is the default for the Boss-Fight dial sheet and for proof-carrying repetition accounting.

\section{A committee-hit calculation, not an observation channel}\label{sec:pobs}
If a committee of size $s$ is sampled uniformly without replacement from a fixed population of $N$ nodes containing exactly $\rho N$ corrupt nodes, then the probability that the committee contains at least one corrupt node is
\[
 p_{\mathsf{hit}} = 1-\frac{\binom{(1-\rho)N}{s}}{\binom{N}{s}}.
\]
The approximation $1-(1-\rho)^s$ replaces finite-population sampling by independent draws. Both expressions are calculations for the named sampling model only. They do not prove that a corrupt member sees the complete declared transcript, that the hit event is independent of the secret and revealed transcript, or that committees are independent across contacts or rounds. Consequently $p_{\mathsf{hit}}$ cannot be inserted into an erasure-attenuation or dummy-width formula without a separate channel witness. Eclipse behavior, route capture, adaptive retries, and target-dependent sampling may also invalidate the sampling model itself.

\section*{How to use this note}
When a paper needs ``threat model'' boilerplate, it should:
(i) cite this note for the TW checklist; (ii) cite the trace-interface note for the tier/exposure normal form; and
(iii) state only the field instantiations (secret, tier, unit, window, universe, channel/witness, active power, and background model) needed for the local theorem. A scalar observation row may accompany those fields as telemetry, never as their substitute.

\begin{thebibliography}{99}\small
\bibitem{series:artifactpolicy}
Anonymous.
\newblock Artifact and Disclosure Policy for the One-True-Archive (Source-Only Public Release).
\newblock Series synthesis note (Synthesis~6), 2026. (This archive.)

\bibitem{series:traceifaces}
Anonymous.
\newblock Trace Interfaces for Anonymous DHTs: Observation Tiers, Committee-Contact Traces, and Exposure Normal Forms.
\newblock Series synthesis note, 2026. (This archive.)

\bibitem{series:deployedsurfaces}
Anonymous.
\newblock Deployed Observation Surfaces for Anonymous DHTs: Control-plane knobs in IPFS/libp2p delegated routing and OHTTP evolution.
\newblock Series synthesis note (Synthesis~30), 2026. (This archive.)


\bibitem{series:receiptitems}
Anonymous.
\newblock Receipt Line Items for Anonymity Claims: A Minimal Tuple Schema for Published Budgets.
\newblock Series synthesis note (Synthesis~12), The-One-True-Archive (2026).

\bibitem{series:workedexample}
Anonymous.
\newblock Worked Example: Instantiating a Tiered Reader-Privacy Receipt.
\newblock Series synthesis note (Synthesis~17), The-One-True-Archive (2026).

\bibitem{series:primarysepmanifest}
Anonymous.
\newblock Primary-Path Separation Manifests for Anonymous-DHT Receipts: Guard Objects for Non-Collusion Assumptions.
\newblock Series synthesis note (Synthesis~22), 2026. (This archive.)

\bibitem{series:probsep}
Anonymous.
\newblock Probabilistic Separation Budgets: Separation Failure as a Receipt-Grade $(\varepsilon,\delta)$ Adjunct.
\newblock Series synthesis note (Synthesis~23), 2026. (This archive.)

\bibitem{DouceurSybil2002}
J.~R. Douceur.
\newblock The sybil attack.
\newblock In \emph{IPTPS}, 2002.

\bibitem{SinghEclipse2004}
A.~Singh, T.~Ngan, and P.~Druschel.
\newblock Eclipse attacks on overlay networks: threats and defenses.
\newblock In \emph{IEEE INFOCOM}, 2006. (Initial versions circulated 2004.)

\bibitem{WangMittalBorisovCCS2010}
Q.~Wang, P.~Mittal, and N.~Borisov.
\newblock In search of an anonymous and secure lookup: attacks on structured peer-to-peer anonymous communication systems.
\newblock In \emph{ACM CCS}, 2010.

\bibitem{MittalBorisov2012InfoLeaks}
P.~Mittal and N.~Borisov.
\newblock Information leaks in structured peer-to-peer anonymous communication systems.
\newblock \emph{ACM TISSEC}, 15(1), 2012.
\end{thebibliography}

\end{document}
